% Editorial addition; distinct from the archived introduction.
\section{Mikhail Kravchuk (1892--1942): life and mathematical legacy}
\label{sec:kravchuk-biography}



Mikhail Kravchuk, also known as Mykhailo Pylypovych Kravchuk, was a
Ukrainian mathematician whose work ranged across algebra, differential
and integral equations, probability, and mathematical statistics. He was
born on 27 September 1892 in Chovnitsy, in Volyn, in what is now Ukraine,
and died on 9 March 1942 in a Soviet forced-labour camp in Kolyma.
The biographical outline here draws on the MacTutor History of Mathematics
Archive and the institutional biography maintained by the Institute of
Mathematics of the National Academy of Sciences of Ukraine.\footnote{MacTutor
History of Mathematics Archive, University of St Andrews,
``Mykhailo Pilipovich Krawtchouk,''
\url{https://mathshistory.st-andrews.ac.uk/Biographies/Krawtchouk/};
Institute of Mathematics, National Academy of Sciences of Ukraine,
``Kravchuk Mikhail Pilipovich (1892--1942),''
\url{https://www.imath.kiev.ua/famous?n=kravchuk&lang=en}.
These are secondary biographical accounts, not archival documents.}

\subsection*{Education and scientific career}
Kravchuk's family moved to Lutsk in 1901. He completed his schooling there
with a gold medal in 1910 and entered the University of Kyiv, then
St Vladimir University, graduating in 1914. His early research included
commuting matrices. After the upheavals of the First World War and the
subsequent revolutionary and civil-war years, he combined university
teaching with work in schools. He defended his doctoral dissertation in
1924 and received the title of professor in 1925. His work on matrices,
bilinear forms, and linear transformations was an important part of this
early development.

In 1929 he was elected a full member of the Ukrainian Academy of Sciences.
He taught at the Kyiv Polytechnic Institute and served as head of its
mathematics chair. The Institute of Mathematics identifies him as one of
its founders. His activity was not confined to research: he helped develop
Ukrainian-language mathematical education and terminology, wrote teaching
materials, and supported the training of younger mathematicians.

\subsection*{The polynomials and the scope of his work}
The orthogonal polynomials bearing his name were introduced in 1929.
They are associated with the binomial distribution, providing the
historical starting point for the finite orthogonality and transform
identities developed in this Encyclopedia. Their subsequent uses in
coding theory, finite transforms, and algebraic constructions belong to
the later development of the subject; they should not all be attributed
to Kravchuk himself.

Although these polynomials are now his most familiar mathematical legacy,
his research was considerably broader. It included matrix theory, the
method of moments, differential and integral equations, probability,
statistics, and questions in mathematical education. This wider setting
is particularly relevant here: the interplay between matrices,
orthogonality, and operator methods is not merely an accidental grouping
of topics under one name.

\subsection*{Persecution, death, and rehabilitation}
Kravchuk was arrested in 1938 during Stalin's purges and sentenced on
political charges to twenty years of imprisonment. He was sent to Kolyma,
where he died in 1942, aged forty-nine. MacTutor describes coercion and
threats against his family during the proceedings. He was posthumously
rehabilitated in 1956. His membership in the Ukrainian Academy of Sciences
was restored in 1992. These events are essential to an account of his
career: its abrupt end was a consequence of political repression, not a
natural conclusion to his scientific work.\footnote{For the account of
his persecution, rehabilitation, and restoration of Academy membership,
see the MacTutor biography cited above.}

\subsection*{Names and conventions in this Encyclopedia}
The spellings \emph{Kravchuk}, \emph{Krawtchouk}, and \emph{Krawtschuk}
reflect different transliterations. MacTutor notes that \emph{Krawtchouk}
was the spelling used in his French-language papers. We retain
\emph{Kravchuk} in the Encyclopedia's title and historical discussion,
while respecting \emph{Krawtchouk} in established mathematical terminology,
source titles, and quotations. The first names \emph{Mikhail} and
\emph{Mykhailo} likewise occur in the biographical literature. These
variants refer to the same mathematician.

\medskip
\noindent\emph{Editorial provenance.} This sourced biographical section
is a later addition to the Encyclopedia. It is not part of the archived
2001 introduction, whose wording and attribution are retained separately.
